\chapter{数学公式与符号的极简美学}

\section{行内、行间与括号}
我把变量和短条件放在句子里，例如 \(x>0\)、\(f(x)=\ln x\)。
需要停下来读的结论则单独成行：
\[
  \int_0^1x^2\,\mathrm{d}x=\frac13.
\]
需要引用时，使用自动编号：
\begin{equation}\label{eq:two:ftc}
  \int_a^b f(x)\,\mathrm{d}x=F(b)-F(a).
\end{equation}
这里假设 \(f\) 在 \([a,b]\) 上连续，\(F\) 在 \([a,b]\) 上连续，
并且在 \((a,b)\) 内满足 \(F'=f\)。
式~\eqref{eq:two:ftc} 的引用通常要编译两遍才稳定。

在源码里，\verb|\(| 与 \verb|\)| 把短公式留在句子里，
\verb|\[| 与 \verb|\]| 让公式另起一行且不编号。
这些定界符本身不印出来。要在公式里印出集合花括号，用 \verb|\{| 和 \verb|\}|。
\[
  (a+b),\quad [a,b],\quad \{x\in\mathbb R\mid x>0\},
  \quad \left(\frac{a+b}{c}\right)^2,
  \quad \left.\frac{x^3}{3}\right|_0^1.
\]
\verb|\left| 与 \verb|\right| 必须配对。单边竖线用点补上不可见的一侧。
别把一对自动伸缩括号直接拆到 align 的两行。

\section{align：把等号对齐}
我用高数笔记中的分部积分题演示连续推导：
\begin{align}
  \int x\mathrm{e}^x\,\mathrm{d}x
    &=x\mathrm{e}^x-\int\mathrm{e}^x\,\mathrm{d}x\notag\\
    &=x\mathrm{e}^x-\mathrm{e}^x+C\notag\\
    &=(x-1)\mathrm{e}^x+C.
\end{align}
\verb|&| 放在等号前，两个反斜杠结束一行。\verb|\notag| 取消该行编号。
align 自己就是行间环境，不要再套一层数学定界符。
整组只要一个编号时，可以这样写：
\begin{equation}
  \begin{aligned}
    (x+1)^2&=x^2+2x+1,\\
    (x-1)^2&=x^2-2x+1.
  \end{aligned}
\end{equation}
aligned 是内部结构，外层 equation 负责整组编号。

\section{矩阵与分段函数}
线代中的方程组可以写成
\[
  \begin{pmatrix}1&1\\2&-1\end{pmatrix}
  \begin{pmatrix}x\\y\end{pmatrix}
  =\begin{pmatrix}3\\0\end{pmatrix},\qquad
  \begin{pmatrix}x\\y\end{pmatrix}=\begin{pmatrix}1\\2\end{pmatrix}.
\]
pmatrix 自带圆括号，\verb|&| 分列，两个反斜杠换行。
先数每行的列数，再怀疑编译器，顺序别反了。

cases 的左列写结果，右列写条件：
\[
  g(x)=\begin{cases}-x,&x<0,\\x,&x\geq0.\end{cases}
\]
概率论里的正部截断也可以直接排成分段形式。
设 \(X\sim U(-1,1)\)，\(Y=\max\{X,0\}\)，则
\[
  F_Y(t)=\begin{cases}
    0,&t<0,\\[3pt]
    \dfrac{t+1}{2},&0\leq t<1,\\[3pt]
    1,&t\geq1.
  \end{cases}
\]
零点属于中间一段，\(F_Y(0)=1/2\)，不能把这个概率跳跃删掉。

\section{向量与希腊字母}
在这套数学字体下，mathbf 常用于粗直立拉丁字母，bm 可处理粗斜体与希腊字母：
\[
  x,\quad\mathbf{x},\quad\bm{x},\qquad
  \alpha,\quad\bm{\alpha},\qquad\mathbf{A}\bm{x}=\bm{b}.
\]
我只加粗向量主体，例如 \(\bm{x}_i\)，下标通常保持普通字重。
函数名使用运算符命令，例如 \(\sin x\)、\(\ln x\)。
少一个反斜杠，有时不报错，只是把函数排成几个散步的字母。

\section{高频符号速查}
中列给出数学模式内的输入，右列显示效果。
\begin{center}
\begin{tabular}{lll}
\toprule
用途 & 输入 & 效果 \\
\midrule
分式 & \verb|\frac{a}{b}| & \(\frac ab\) \\
根式 & \verb|\sqrt{x}| & \(\sqrt{x}\) \\
导数 & \verb|f'(x)| & \(f'(x)\) \\
偏导 & \verb|\frac{\partial f}{\partial x}| & \(\frac{\partial f}{\partial x}\) \\
积分 & \verb|\int_a^b f(x)\,\mathrm{d}x| & \(\int_a^b f(x)\,\mathrm{d}x\) \\
极限 & \verb|\lim_{x\to0}f(x)| & \(\lim_{x\to0}f(x)\) \\
求和 & \verb|\sum_{i=1}^n x_i| & \(\sum_{i=1}^n x_i\) \\
\bottomrule
\end{tabular}
\end{center}
\begin{center}
\begin{tabular}{lll}
\toprule
用途 & 输入 & 效果 \\
\midrule
希腊字母 & \verb|\alpha,\beta,\gamma| & \(\alpha,\beta,\gamma\) \\
常用参数 & \verb|\theta,\lambda,\sigma| & \(\theta,\lambda,\sigma\) \\
大小写 & \verb|\delta,\Delta,\Omega| & \(\delta,\Delta,\Omega\) \\
数集 & \verb|\mathbb{R},\mathbb{N}| & \(\mathbb{R},\mathbb{N}\) \\
属于 & \verb|\in,\notin| & \(\in,\notin\) \\
包含 & \verb|\subseteq,\subsetneq| & \(\subseteq,\subsetneq\) \\
交并 & \verb|A\cap B,A\cup B| & \(A\cap B,A\cup B\) \\
空集 & \verb|\varnothing| & \(\varnothing\) \\
\bottomrule
\end{tabular}
\end{center}
符号只记得形状时，我会用 \href{https://detexify.kirelabs.org/}{Detexify}
画出来查命令，再检查是否需要额外宏包。

\section*{本章完整源码}
\sourceattachment{clean_chapters/ch02.tex}
\VerbatimInput{clean_chapters/ch02.tex}
